\section{Implementation Architecture and Design}

Apollo-RTOS is a real-time operating system inspired by the AGC, designed for the
\texttt{micro:bit} v1 device. It follows a monolithic kernel architecture and is built
around a cooperative scheduler with preemptive capabilities. 

This section provides an overview of the core modules in Apollo-RTOS and their
interconnections, followed by a discussion of the codebase structure.

\subsection{Overview}

The major modules of Apollo-RTOS are:

\paragraph*{Scheduler} The core of the OS, responsible for:
\begin{itemize}
    \item Managing process execution, including creation, termination, priority
        changes, pausing, and resuming.
    \item Handling context switching between processes.
\end{itemize}

\paragraph*{Waitlist} Implements a deadline-driven task scheduler, similar to cron
jobs, enabling time-sensitive operations to be executed at precise intervals.

\paragraph*{Signals and Mutexes} Enable low-overhead inter-process communication (IPC)
and synchronization mechanisms among the processes.

\paragraph*{Memory Manager} Handles all dynamic memory allocation. It implements a
freelist-based \texttt{malloc} for allocating and deallocating runtime memory and
manages process stack allocations.

\paragraph*{File System} Provides persistent storage, implemented using an external
FRAM over I2C. An optional Flash-based file system can also be enabled.

\paragraph*{Debugging and Error Handling} Provides essential debugging utilities,
including:
\begin{itemize}
    \item \texttt{debug}: Outputs diagnostic information.
    \item \texttt{assert}: Handles runtime error detection.
\end{itemize}

\paragraph*{Recovery} Ensures system stability by protecting critical processes and
tasks from various resets and failures.

\paragraph*{Drivers} Interface with the hardware components of the \texttt{micro:bit}.
The key drivers include:
\begin{itemize}
    \item \textbf{Timer}: Apollo-RTOS utilizes two 16-bit timers—\texttt{Timer1} as a
        system clock and \texttt{Timer2} for profiling.
    \item \textbf{Serial}: Manages UART communication.
    \item \textbf{I2C}: Handles communication with the I2C master on the
        \texttt{micro:bit}.
    \item \textbf{Display}: Controls the \(5 \times 5\) LED matrix on the board.
\end{itemize}

\paragraph*{Shell} Implements a Unix-like command-line interface that allows users to
manage and debug processes over UART. It maintains a list of available commands mapped
to various processes.


